% TOP_PROBLEM: {"id": 5300043, "title": "Lebesgue ergodicity on a spherical Julia set", "category_id": 11, "category": {"id": 11, "name": "dynamical_systems", "display_name": "Dynamical Systems", "description": "Problems about long-term behavior of deterministic systems, Hamiltonian dynamics, and periodic orbits.", "slug": "dynamical-systems", "order_index": 11, "created_at": "2026-07-31T15:26:25.671Z"}, "status": "open", "problem_number": "AMR-052-0043", "set_id": 13}
% FIRST_POSTED: 2026-10-02
% ENTRY_KIND: statement_only
% UnsolvedMath ID 5300043; source catalogue code AMR-052-0043
% !TeX program = lualatex
% Definition and sources only; this file is not an LLM solution attempt.
\documentclass[11pt,a4paper]{article}
\usepackage[margin=25mm]{geometry}
\usepackage{fontspec}
\setmainfont{lmroman10-regular.otf}[BoldFont=lmroman10-bold.otf,
 ItalicFont=lmroman10-italic.otf,BoldItalicFont=lmroman10-bolditalic.otf]
\usepackage{amsmath,amssymb,mathtools}
\usepackage{unicode-math}
\setmathfont{latinmodern-math.otf}
\AtBeginDocument{\let\setminus\smallsetminus}
\usepackage{xurl,hyperref}
\hypersetup{colorlinks=true,urlcolor=blue}
\setcounter{secnumdepth}{0}
\setlength{\parindent}{0pt}
\setlength{\parskip}{5pt}
\setlength{\emergencystretch}{3em}
\begin{document}
\section*{Lebesgue ergodicity on a spherical Julia set}\label{5300043}
{\small UnsolvedMath ID 5300043 · AMR-052-0043 · Dynamical Systems · AMR Open Problem Lists}
\subsection{Definitions and mathematical statement}
Let $R:\widehat{\mathbb C}\to\widehat{\mathbb C}$ be a rational map of
degree $d\ge2$, and assume its Julia set is the whole Riemann sphere.
Equivalently there is no nonempty open set on which the iterates form
a normal family. Let $m$ be normalized spherical area measure.
Although $R$ need not preserve $m$, it is nonsingular: inverse images
of area-zero sets have area zero.

The first question is whether $R$ is ergodic for this measure class,
meaning that every measurable $E$ with
$m(R^{-1}E\mathbin{\triangle} E)=0$ has $m(E)=0$ or $m(E)=1$.
This definition does not assert that spherical area itself is an
invariant probability measure.

If universal ergodicity fails, the source asks the weaker bound of
at most $2d-2$ ergodic components. Precisely, can the sphere, modulo
an area-null set, be partitioned into $r\le2d-2$ invariant measurable
sets of positive area, on each of which the restricted nonsingular
action is ergodic? The parts are understood modulo null sets, and
invariance is under full inverse image in that same sense.

The hypothesis excludes attracting basins, Siegel disks, and all other
Fatou components. Ergodicity for a different preferred invariant
measure, such as the measure of maximal entropy, does not answer
the area-class question. The first assertion is stronger than the
finite-component alternative, and either requires control of all
measurable invariant subsets, not only closed invariant sets.
\subsection{Short English statement}
Must a rational map with spherical Julia set be ergodic for the area measure class, or have at most $2d-2$ components?
\subsection{Sources}
\begin{itemize}
\item[S1] ULAM.AI and the UnsolvedMath Contributors, supplied problems.json, record 5300043 (AMR-052-0043), AMR Open Problem Lists. Catalogue identity and source transcription; CC BY 4.0.
\url{https://huggingface.co/datasets/ulamai/UnsolvedMath}.
\item[S2] Stony Brook - Open Problems in Dynamical Systems, 1992, Lyubich P3. Original source indexed by AMR-052-0043.
\url{https://www.math.stonybrook.edu/open-problems-dynamical-systems}.
\item[S3] Lyubich, Problems on measurable dynamics, Problem3, in Bielefeld–Lyubich (1992).
\url{https://www.math.stonybrook.edu/preprints/ims92-7.pdf}.
\end{itemize}
\end{document}
